%% kapa: exam=1 round=33 year=2022 period=2 form=A subject=accounting subject_ko=회계학 seq=28 number=68
%% source: Q-Net 감정평가사 기출문제 (cst00309) artlSeq=5209997 file=raw/1008/5209997/2022년도 제33회 감정평가사 1차시험 2교시(감정평가 관계법규, 회계학).pdf
%% conversion: kordoc 4.24.51 -> tools/gen1.py (auto); review=proofread-sonnet; flags=none
\begin{question}{68}
(주)감평은 20×1년 초 ￦20,000에 기계장치(내용연수 5년, 잔존가치 ￦0, 정액법 감가상각)를 취득하여 사용하고 있다. (주)감평은 동 기계장치에 대해 취득연도부터 재평가모형을 적용하고 있으며, 처분부대원가가 무시할 수 없을 정도로 상당하여 손상회계를 적용하고 있다. 공정가치와 회수가능액이 다음과 같을 경우, 20×2년도에 인식할 손상차손 또는 손상차손환입액은? (단, 기계장치를 사용함에 따라 재평가잉여금의 일부를 이익잉여금으로 대체하지 않는다.)
\par{\small\centering\begin{tabular}{|c|r|r|}\hline
구 분 & 20×1년 말 & 20×2년 말 \\ \hline
공정가치 & ￦18,000 & ￦12,000 \\
회수가능액 & 19,500 & 11,000 \\ \hline
\end{tabular}\par}
\begin{choices}
  \item ￦0
  \item 손상차손 ￦500
  \item 손상차손 ￦1,000
  \item 손상차손환입 ￦500
  \item 손상차손환입 ￦1,000
\end{choices}
\end{question}
\answer{2,3}
\answernote{복수 정답(정답표 원문: 2, 3)}
